\section{Introduction}
\label{sec:intro}

The strictly convex quadratic program
\begin{equation}
  \min_{x \in \R^n} \; \tfrac12 x\T G x - a\T x
  \qquad \text{subject to} \qquad C\T x \ge b,
  \label{eq:qp}
\end{equation}
with $G \in \R^{n \times n}$ symmetric positive definite, $C \in \R^{n \times m}$
and the first $m_{\mathrm{eq}}$ of the $m$ constraints understood as equalities,
is solved a great many times a day. Mean-variance portfolio
selection~\cite{markowitz1952} is exactly~\eqref{eq:qp}; so is every step of a
sequential quadratic programming method and every horizon of a linear model
predictive controller. For dense instances of small to moderate size --- $n$ from
a handful to a few thousand --- the method of choice has for forty years been the
dual active-set algorithm of Goldfarb and Idnani~\cite{goldfarb1982,goldfarb1983},
whose distinguishing property is that it needs no phase-one problem: the
unconstrained minimiser $G^{-1}a$ is dual feasible for free, so the iteration can
start there and drive primal infeasibility to zero.

This note is a derivation of that method in the form the \texttt{cvx-quadprog}
package~\cite{cvxquadprog} implements it, together with the two extensions that
package adds. It is written because the companion paper accompanying the
package~\cite{schmelzer2026quadprog} is about something else. That paper asks
which of the reference implementation's decisions are essential to the algorithm
and which are artefacts of Fortran, and answers with measurements; it states the
mathematics compactly and moves on. What follows is the part it moves past: every
identity the solver depends on, derived, with the degenerate cases named. No
measurement is repeated here: where a claim below rests on one, it is cited to
that paper rather than restated as if this one had established it.

\subsection{Why the dual method}

A primal active-set method maintains a feasible iterate and works towards
stationarity. It must therefore begin at a feasible point, and finding one is a
linear program in its own right --- the phase-one problem. The dual method
reverses the roles. It maintains stationarity and the sign conditions on the
multipliers throughout, and works towards feasibility. Its starting point is
free, because with no constraint active the stationarity condition reads
$Gx = a$ and the multiplier vector is empty, so $x = G^{-1}a$ satisfies
everything the method asks of an iterate. One Cholesky factorisation and one
triangular solve replace the phase-one problem entirely.

The price is that the iterate is infeasible until the last step, so the method
cannot be stopped early for an approximate answer. In exchange it terminates at
an exactly feasible point of the active-set lattice rather than approaching one
asymptotically as an interior-point method does, and it warm-starts almost
perfectly, which is why parametric solvers are built on active-set
methods~\cite{nocedal2006}. Section~\ref{sec:sweep} exploits that.

\subsection{What is derived here}

Section~\ref{sec:problem} fixes notation, records the optimality conditions, and
establishes the two facts the rest of the note leans on: that~\eqref{eq:qp} has at
most one solution, and that its Karush--Kuhn--Tucker conditions are sufficient
and not merely necessary. The second is what makes a certificate possible, and
both of the extensions in Sections~\ref{sec:pdas} and~\ref{sec:sweep} are built
on certificates.

Section~\ref{sec:factorisation} introduces the single matrix the solver carries
and the two invariants that define it. This is the heart of the method and the
part most easily mistaken for an implementation detail. The matrix is not a
factor of $G$, nor an orthogonal factor of the constraint normals; it is the
product of both, and the reason to fold them together is that it makes the search
directions matrix--vector products of fixed cost rather than solves whose cost
grows with the active set. We give the geometric reading --- the columns are a
$G$-orthonormal basis of $\R^n$ split so that the trailing block spans the
working set's feasible subspace --- and then show, in
Proposition~\ref{prop:invariance}, that the two quantities the iteration consumes
depend on the invariants alone and not on which of the many matrices satisfying
them is held. That is the licence for the implementation's choice of orthogonal
reduction.

Section~\ref{sec:iteration} derives the iteration: the two search directions and
their closed forms, the fact that the multipliers move affinely along $-r$ while
the entering multiplier rises with the step, the two competing step lengths, and
the exact objective increment. Section~\ref{sec:termination} collects what can be
proved about stopping --- strict monotonicity, termination of the inner loop in at
most one pass per active constraint, and the Farkas certificate that turns a stuck
iteration into a proof of infeasibility. Section~\ref{sec:updates} derives the two
factorisation updates.

Sections~\ref{sec:pdas} and~\ref{sec:sweep} treat the package's two departures
from the classical method, and Section~\ref{sec:arithmetic} the finite-precision
questions: how the constraint-selection rule is made invariant to how each
constraint happens to be scaled, and why a violation the size of the rounding in
the iterate must be set aside rather than either enforced or taken as evidence of
infeasibility.

\subsection{Notation}

Vectors are columns. For an index set $\Active \subseteq \{1,\dots,m\}$ of size
$k$ we write $A = C_{\cdot\Active} \in \R^{n \times k}$ for the matrix of the
corresponding constraint normals, $b_\Active$ for the corresponding right-hand
sides, and $c_i$ for the $i$-th column of $C$. Throughout, $\Active$ is the
solver's \emph{working set}: the constraints it is currently holding as
equalities. At a solution the working set is a subset of the constraints active
there, and the two coincide except in the degenerate case treated in
Remark~\ref{rem:degeneracy}. The entering constraint's normal is $n_*$ and its
right-hand side $b_*$. We reserve $u$ for multipliers of working-set constraints,
$u_*$ for the entering constraint's own multiplier, and $\lambda$ for a full
length-$m$ multiplier vector with zeros off the working set.
